\section{Observation frame and the velocity in the denominator}
\label{sec:frame-mass}

The distinction between the entire spatial plateau and the complete
velocity is essential.  In this section write
$M_{F,L}=M_L^{\mathrm{plat}}$ for
\eqref{eq:parabolic-physical-payment}, and retain the exact initial
coefficients, cutoffs, and degree sequences of
Theorem~\ref{thm:parabolic-shear} unless stated otherwise.
The following consequences concern the normalization of that explicit
family.  They are not general local-energy estimates.

\begin{proposition}[A change of observation frame]\label{prop:galilean-frame}
For $u(t,x)=(0,B,F(t,x_2))$, define $y=x-Bt e_2$ and
\[
 \widetilde u(t,y)=u(t,y+Bt e_2)-B e_2=(0,0,H(t,y_2)).
\]
Then $H_t=H_{y_2y_2}$.  For any smooth radial cutoff $\phi(y)$
supported in a Euclidean coordinate chart of the torus,
\[
 \frac12\int_{\mathbb T^3}|\widetilde u|^2
                     \widetilde u\cdot\nabla\phi\,dy=0.
\]
\end{proposition}
\begin{proof}
The definition gives $H(t,y_2)=F(t,y_2+Bt)$, so
$H_t=F_t+BF_{x_2}=F_{x_2x_2}$.
The flux integrand is $H(t,y_2)^3\partial_{y_3}\phi/2$.
Its integral in $y_3$ vanishes by the fundamental theorem of calculus
and the support of $\phi$.  Periodicity gives the same conclusion.
\end{proof}

A cutoff fixed in $y$ corresponds to a translating cutoff in the
original coordinates.  Moreover the kinetic energy changes when the
constant velocity is subtracted.  Thus this zero flux and the positive
flux in Theorem~\ref{thm:parabolic-shear} are different observations;
the latter is not Galilean invariant.  It measures constant transport
through a fixed window.  The construction does not introduce a
vortex-stretching or energy-cascade mechanism.

\begin{proposition}[Complete velocity mass for the stated amplitude]
\label{prop:full-velocity-mass}
Define the positive full-velocity mass and the spacetime volume by
\[
 M_{u,L}=\int_0^{4R^2}\int_{\mathcal P_L}
                  (B^2+F(t,x_2)^2)^{3/2}\,dx\,dt,
 \qquad V_L=4R^2|\mathcal P_L|.
\]
Along every degree sequence \eqref{eq:edge-degree-regime},
\begin{align}
 V_L&=32\pi\delta_0a^2R^5
                   \left(1+\frac{\delta_0^2}{3a^2}\right),
                   \label{eq:complete-volume}\\
 M_{u,L}&=(1+o(1))\,32\pi\delta_0\beta^3
                       \frac{a^5}{A^3}R^2,\label{eq:complete-mass}\\
 M_{F,L}&=R^5\exp(o(L^2)),\qquad
 \mathcal T_L=R^3\exp(o(L^2)).\label{eq:component-scales}
\end{align}
The final expression uses the eventual positivity of $\mathcal T_L$.
In particular
\begin{equation}\label{eq:complete-velocity-rate}
 R^{1/3}\frac{\mathcal T_L}{M_{u,L}^{2/3}}
   =R^2\exp(o(L^2))\longrightarrow0,
 \qquad
 \lim_{L\to\infty}\frac1{L^2}
 \log\left(R^{1/3}\frac{\mathcal T_L}{M_{u,L}^{2/3}}\right)
   =-\frac9{20000}.
\end{equation}
All limits refer to the amplitudes fixed in
\eqref{eq:parabolic-fourier-coefficients}; they are not uniform over
arbitrary rescalings of those amplitudes.
\end{proposition}
\begin{proof}
Subtracting the volumes of the two balls gives
\eqref{eq:complete-volume} exactly.  Since $a,A\asymp L$ and
$\mu=o(A)$, the global majorant and dynamic transfer give
\[
 \sup_{[0,4R^2]\times\mathcal P_L}|F|
 \le C Q_{m,A}(4,c_p+4\gamma_L)+o(1)
 =\exp(o(L^2)),
 \qquad \gamma_L=\beta a/A.
\]
Indeed the heat-edge theorem gives, for every bounded coordinate $c$,
$\log Q_{m,A}(4,c)=F_4\mu^2+o(\mu^2)+O(\mu)$ uniformly in $c$.
On the other hand $|B|=\beta a/(AR)$, and hence
$\sup|F|/|B|\to0$.  Uniformly on the entire integration region,
$(B^2+F^2)^{3/2}=|B|^3(1+o(1))$.
Multiplication by the exact volume proves
\eqref{eq:complete-mass}.

For the remaining two scales, use the positive lower and upper bounds
\eqref{eq:parabolic-S-bounds}, \eqref{eq:parabolic-K-upper}, and
\eqref{eq:parabolic-K-lower}, followed by
\eqref{eq:parabolic-dimensional-identities}.
Every $Q(4,c)$ in those bounds has logarithm $o(L^2)$;
$a,A,\beta$, and $h_\mu=(1+\mu^2)^{-1}$ contribute logarithms
$O(\log A)=o(L^2)$.  Both sides therefore give
\eqref{eq:component-scales}.  This step uses a lower bound on the
whole signed time integral, not an absolute value of a cap contribution.

The denominator $M_{F,L}$ is positive for every nonzero member of the
family: a nonzero finite Fourier sum cannot vanish on a spatial
interval, and the plateau has positive cross-sectional area on an
interval.  The time evolution preserves its nonzero Fourier
coefficients.  The positivity of $\mathcal T_L$ is only eventual and
comes from paired late-time caps and absorption of the earlier signed
contribution in Theorem~\ref{thm:parabolic-shear}.
Thus the logarithm in \eqref{eq:complete-velocity-rate} is defined
eventually along each sequence.  Equations
\eqref{eq:complete-mass}--\eqref{eq:component-scales} give the net
power $R^{1/3+3-4/3}=R^2$.  Finally
$2L^{-2}\log R=-9/20000$.
\end{proof}

The qualifier about amplitude cannot be removed.  The next elementary
observation makes that dependence explicit.

\begin{proposition}[Amplitude rescaling]\label{prop:amplitude-scaling}
Fix one shear, a finite observation region of spacetime volume $V>0$,
and the same numerator and denominator conventions as above.  Assume
$M_F>0$ and $\mathcal T\ge0$.  For $\lambda>0$ replace its initial
profile by $\lambda F(0,\cdot)$ and keep $B$ fixed.  The resulting
solution has profile $\lambda F(t,\cdot)$ throughout the same interval.
Set
\[
 Z_F=R^{1/3}\frac{\mathcal T}{M_F^{2/3}},\qquad
 Z_u(\lambda)=R^{1/3}
  \frac{\lambda^2\mathcal T}
  {\left[\int(B^2+\lambda^2F^2)^{3/2}\right]^{2/3}}.
\]
Then $Z_u(\lambda)$ is nondecreasing and
\begin{equation}\label{eq:amplitude-sandwich}
 \frac{Z_F}{1+(B/\lambda)^2(V/M_F)^{2/3}}
 \le Z_u(\lambda)\le Z_F,
 \qquad \lim_{\lambda\to\infty}Z_u(\lambda)=Z_F.
\end{equation}
For the family of Proposition~\ref{prop:full-velocity-mass}, the
explicit choice $\lambda_L=R^{-2}$ gives
$Z_u(\lambda_L)/Z_F\to1$.  Consequently that rescaled family has
$Z_u(\lambda_L)\to\infty$ if $\mu/\log A\to\infty$.
\end{proposition}
\begin{proof}
The scalar advection--heat equation is linear, and its transport flux
is quadratic in $F$.  After dividing the full mass by $\lambda^3$,
its two-thirds power is
\[
 \left\|F^2+(B/\lambda)^2\right\|_{L^{3/2}}.
\]
It decreases with $\lambda$ and is bounded below by
$\|F^2\|_{L^{3/2}}=M_F^{2/3}$.  Minkowski's inequality bounds it
above by $M_F^{2/3}+(B/\lambda)^2V^{2/3}$.  This proves the
inequalities, monotonicity, and limit, including $B=0$, when equality
holds throughout.  Dominated convergence also gives the limit directly,
with dominator $(F^2+B^2)^{3/2}$ for $\lambda\ge1$.
In the specified family,
\[
 (B/\lambda_L)^2(V_L/M_{F,L})^{2/3}
       =R^2\exp(o(L^2))\longrightarrow0.
\]
The asserted divergence now follows from
Theorem~\ref{thm:parabolic-shear}.  Each rescaled member still uses
one initial datum over the entire time interval.
\end{proof}

This rescaling retains the very large constant drift and does not
produce uniformly bounded initial energy.  It shows why
\eqref{eq:complete-velocity-rate} cannot be read as an impossibility
theorem for every amplitude normalization.  It is also not a
counterexample to a general local-energy estimate: the component
estimate in this paper retains its stated geometric and frequency
factors, and rescaling does not remove them.

\begin{proposition}[Bounded drift for the same initial shape]
\label{prop:bounded-drift-shape}
Keep $a,A,R,m$ and the initial coefficients
\eqref{eq:parabolic-fourier-coefficients}, but replace the drift by any
$B_L$ with $|B_L|\le B_*<\infty$.  Denote the corresponding flux and
component mass by $\mathcal T_L^{(B_L)}$ and $M_{F,L}^{(B_L)}$.
There is a constant $C$, depending on the fixed cutoffs and $B_*$,
such that eventually along every allowed degree sequence
\begin{equation}\label{eq:bounded-drift-decay}
 R^{1/3}\frac{|\mathcal T_L^{(B_L)}|}
 {(M_{F,L}^{(B_L)})^{2/3}}
 \le C |B_L|R a^{1/3}(1+\mu^2)^{2/3}e^{C\mu}
 \longrightarrow0.
\end{equation}
The same upper bound holds after arbitrary positive amplitude
rescaling and replacement of the component mass by the complete
velocity mass.  If $B_L=0$, the flux is exactly zero.
\end{proposition}
\begin{proof}
In Proposition~\ref{prop:dynamic-confluence} the real shift becomes
$B_LRs/A$ in the $x$ coordinate.  All other operator, carrier, and
confluence estimates are unchanged; the shift is smaller than the one
already permitted there.  Thus the same absolute transfer bound holds
uniformly for bounded $B_L$.

The exact flux cross section and $\int|W_a|\le C$ give
\[
 |\mathcal T_L^{(B_L)}|
       \le C|B_L|aR^4 Q_{m,A}(4,C_1)^2
\]
for a fixed $C_1>0$.  For the mass lower bound choose the outer
plateau strip on the side whose displacement is $+|B_L|Rs$ in the
edge coordinate.  Parity of the ideal polynomial permits either side.
On its terminal time layer of length $h_\mu$, its edge coordinate is
at least $\delta_0$.  The terminal derivative estimates and the
absolute transfer bound consequently give
\[
 M_{F,L}^{(B_L)}
       \ge caR^5 h_\mu Q_{m,A}(4,\delta_0)^3.
\]
This lower bound restricts the integration only to prove the inequality;
the defined mass still includes the full plateau and time interval.
The heat-edge ratio between $C_1$ and $\delta_0$ is at most $e^{C\mu}$.
Combining the last two bounds gives
\eqref{eq:bounded-drift-decay}.  Its logarithmic upper bound is
$\log R+O(\mu+\log A)$ plus a bounded term, which tends to
$-\infty$ because $\mu=o(A)$ and $A\asymp L$.
The component quotient is unchanged by amplitude rescaling, whereas
the full mass is at least the component mass.  These facts give the
last assertion.  At $B_L=0$ the exact cross-sectional flux formula
vanishes.
\end{proof}

This last proposition concerns the same initial shape and scale.
It neither rules out other profiles with bounded drift nor supplies
matching dynamic upper and lower bounds over the general Fourier class.
